%
% $Log: logicalcon.tex,v $
% Revision 1.4  2002/10/21 15:26:44  mjm
% to ohio, revised
%
% Revision 1.3  2001/08/15 18:07:03  mjm
% APPM preprint 466
%
% Revision 1.2  1999/08/09 19:57:38  mjm
% beta release
%
%

% Notes:
% 1. should we include inverse?

%%%%%%%%%%%%%%%%%%%%%%%%%%%%%%%%%%%%%%%%%%%%%%%%%%%%%%%%%%%%%%%%%%%%%%%%%
\handout{Logical Connectives}

%% reset equation counter
\setcounter{equation}{0}

Mathematics has its own language.  As with any language, effective
communication depends on logically connecting components.  Even the
simplest ``real" mathematical problems require at least a small amount
of reasoning, so it is very important that you develop a feeling for
{\em formal (mathematical) logic}.

Consider, for example, the two sentences ``There are 10 people waiting
for the bus" and ``The bus
is late."  What, if anything, is the logical connection between these two
sentences?  Does one {\em logically imply} the other?
Similarly, the two mathematical statements ``$r^2 + r - 2 = 0$" and
``$r=1$ or $r=-2$" need to be connected, otherwise they are merely two
random statements that convey no useful information.
  Warning: when
mathematicians talk about implication, it means that one thing {\bf must}
be true as a consequence of another; not that it can be true, or
might be true sometimes.  

Words and symbols that tie statements together logically are called
{\em logical connectives}.  They allow you to communicate the reasoning
that has led you to your conclusion.  Possibly the most important of
these is {\em implication} --- the idea that the next statement is a
logical consequence of the previous one.  This concept can be conveyed by
the use of words such as: therefore, hence, and so, thus, since, if $\dots$
then $\dots$, this implies, etc.  In the middle of mathematical calculations, we
can represent these by the implication symbol ($\Rightarrow$).
   For example
\begin{eqnarray}
x + 7y^2 = 3 & \Rightarrow & y = \pm \sqrt{\frac{3-x}{7}} \, ; \label{first} \\
x \in (0,\infty) & \Rightarrow & \cos(x) \in [-1,1]. \label{second}
\end{eqnarray}

\subsubsection*{Converse}
Note that ``statement A $\Rightarrow$ statement B" does {\bf not}
necessarily mean that the {\em logical converse} --- ``statement B
$\Rightarrow$ statement A" --- is also true.  Logical implication is a
matter of cause and effect; the logical converse is simply the reverse
cause-effect situation (which may not be true).
Consider the following, everyday example:
\begin{itemize}
\item ``I am running to class because I am late."
\end{itemize}
It should be clear that the logical converse of these is not true:
\begin{itemize}
\item ``I am late to class because I am running."
\end{itemize}

For examples (\ref{first}) and (\ref{second}) above:
\begin{eqnarray*}
x + 7y^2 = 3 & \Leftarrow & y = \pm \sqrt{\frac{3-x}{7}} \, ; \\
x \in (0,\infty) & \not\Leftarrow & \cos(x) \in [-1,1]. \quad
	\fbox{Since $x<0 \Rightarrow \cos(x) \in [-1,1]$ also.}
\end{eqnarray*}


\subsubsection*{Contrapositive}
We have seen that simply reversing a logical statement can lead to problems.
There is a way, however, to invert an implication so that the inverted
statement is also true.  This is known as the {\em contrapositive}.
Consider the statement ``$A \Rightarrow B$"; this means ``if $A$ is
true, then $B$ is true".  The contrapositive is ``if $B$ is {\bf false},
then $A$ must be false".  
%Recall our earlier example: a power outage causes
%everyone's TV to turn off.  If someone's TV is on ({\em i.e.} it is not
%true that everyone's TV is off), then there cannot be a power outage.

Here is an example of how these concepts fit together:
\begin{itemize}
\item If X is a cat, it is an animal.  \fbox{Accept as true}
\item If X is an animal, it is a cat.  \fbox{FALSE (converse)}
\item If X is not an animal, it is not a cat.  \fbox{TRUE (contrapositive)}
\item If X is not a cat, it is not an animal.\\
	  \fbox{FALSE. This is the
      contrapositive of the second statement, which is false.}
\end{itemize}

\secondpage


\subsubsection*{Equivalence}
When the implication works both ways, we say that the two statements
are {\em equivalent} and we may use the equivalence symbol
($\Leftrightarrow$); in words we may say ``A is equivalent to B" or
``A if and only if B".  If two statements are equivalent, we may use
any of the implication symbols ($\Rightarrow$, $\Leftarrow$, or
$\Leftrightarrow$).  Which connective we use depends on what we are
trying to show.  In (\ref{first}) above, if we are trying to obtain a
formula for $y$, we would probably just use ``$\Rightarrow$", even
though the stronger statement (``$\Leftrightarrow$") is also true.
If, however, we wanted to show that two statements about $x$ and $y$
were equivalent, we would write $x + 7y^2 = 3
\Leftrightarrow y = \pm \sqrt{\frac{3-x}{7}}$.


Careful logic is the heart and soul of mathematics; learn to reason with
watertight arguments and use logical connectives to explain your reasoning
processes.

\subsubsection*{Example:}

{\em What is the greatest amount of water that a right-cylindrical water tank
can hold if there is $100 m^2$ of material from which to construct it?}

\begin{goodexample}
Let the height of the tank be $h$ and the radius be $r$.  \underline{Then}
the volume of the tank is $V(h,r) = \pi hr^2$ and the surface area is
$S(h,r) = 2\pi r^2 + 2\pi rh = 2\pi r(r+h)$.  \underline{Since} we require
$V > 0$, we can assume that $r,h > 0$.  \underline{If} we have $100 m^2$
of material, \underline{then} $S = 100$, \underline{and so}
\begin{eqnarray*}
S & = & 2\pi r(r+h) = 100 \\
 & \Rightarrow & r + h = \frac{100}{2\pi r} \\
 & \Rightarrow & h = \frac{50}{\pi r} - r \\
 & \Rightarrow & V = \pi hr^2 = 50r - \pi r^3.
\end{eqnarray*}
To find the maximum volume, we look for $r$ such that $\frac{dV}{dr} = 0$.
Differentiating with respect to $r$ gives $\frac{dV}{dr} = 50 - 3\pi r^2$.
\underline{Hence},
\begin{eqnarray*}
\frac{dV}{dr} = 0 & \Leftrightarrow & 50 - 3\pi r^2 = 0 \\
    & \Leftrightarrow & r^2 = \frac{50}{3\pi} \\
    & \Leftrightarrow & r = \pm \sqrt{\frac{50}{3\pi}}.
\end{eqnarray*}
\underline{Since} $r>0$, we can use the second derivative test and find 
$\frac{d^2V}{dr^2} = -6\pi r <0$. 
\underline{This implies} that  $r = \sqrt{\frac{50}{3\pi}} \approx 2.3033$
maximizes $V$.
From above, the optimal volume is
\begin{eqnarray*}
V & = & 50r - \pi r^3 \\
  & = & \left( 50 - \pi \left( \frac{50}{3\pi} \right) \right) \sqrt{\frac{50}{3\pi}} \\
  & = & \left( \frac{100}{3} \right) \sqrt{\frac{50}{3\pi}} \approx 76.7765.
\end{eqnarray*}
\underline{Thus} we have found that the greatest amount of water such a tank
can hold is (approximately) $76.7765 m^3$.
\end{goodexample}



